\documentclass[10pt,twocolumn]{article}

\usepackage[margin=0.75in]{geometry}
\usepackage{times}
\usepackage{microtype}
\usepackage{graphicx}
\usepackage{booktabs}
\usepackage{amsmath}
\usepackage{enumitem}
\usepackage{natbib}
\usepackage[hidelinks]{hyperref}
\usepackage{tcolorbox}
\usepackage{xcolor}
\usepackage{tikz}
\usepackage{tabularx}
\urlstyle{same}

\setlength{\columnsep}{0.25in}
\setlength{\parindent}{1em}
\setlength{\parskip}{0pt}
\raggedbottom

% ─────────────────────────────────────────────────────────────────────────────
% Conference first-stage submission: abstract + introduction only.
% The introduction is the detailed ("complex") version and is supplemented with
% the concrete CCA rule so that abstract + introduction are self-contained.
% ─────────────────────────────────────────────────────────────────────────────

\title{\textbf{Cross-Constituency Aggregation for Community Notes}}

\author{
Emrecan Ulu, Jingyao Shi, Carina I. Hausladen\\
University of Konstanz
}

\date{}

\begin{document}

\maketitle

\begin{abstract}
Community Notes is X's crowd-based fact-checking system: contributors write
contextual notes, other contributors rate them, and a bridging algorithm decides
which notes become visible. The algorithm is designed to avoid simple majority
capture by estimating note quality separately from ideological alignment, so a
note should gain support across disagreement rather than only within one side
\citep{wojcik2022birdwatch,xcommunitynotes2026ranking}. This design is powerful,
but it also creates a visibility problem. Because the system works through
sparse and unequal participation, display decisions can depend heavily on the
ratings supplied by a relatively small active core, and many notes remain in
Needs More Ratings when the algorithm does not find enough diverse support
\citep{goyal2026qsmf,razuvayevskaya2025timeliness}. We argue that this is not
only a data-sparsity problem. It is also an aggregation problem: Community Notes
treats cross-group agreement indirectly, through a fitted model, rather than as
an explicit decision rule over constituencies.

We take inspiration from political science and social choice. In systems such
as Switzerland's double-majority referendum rule, important decisions must gain
support not only from a national majority but also from a majority of cantons;
the system treats cross-group support as part of the decision's legitimacy. Our
philosophy is simple: a note should not become visible because one active group
overwhelms the rest; it should become visible when each relevant constituency
has a real chance to support it. We operationalize this idea as
cross-constituency aggregation. We recover behavioral constituencies from
co-rating patterns among 200,000 active raters, compute each note's approval
rate inside each constituency, and combine those rates with a geometric mean
under a minimum-coverage rule. In a 100,000-note analytical slice, Community
Notes displays 6,832 of 44,722 representative-picked notes (15\%). Our rule
identifies a 13,655-note rescue pool, raising potential visibility to 20,487
notes (46\%) before validation. A second-stage validation screen classifies
3,896 candidates as sourced, factual, and rescue-worthy. These results suggest
that useful notes can remain hidden when cross-group support is modeled only
indirectly.
\end{abstract}

\section{Introduction}

Community Notes is X's crowdsourced fact-checking system. When users believe a
post needs context, contributors can write a short note. Other contributors
then rate whether that note is helpful. If the system decides that the note has
enough support, the note is shown under the post. This makes Community Notes a
large public experiment in crowd-based moderation: the platform does not rely
only on professional fact-checkers or automated classifiers, but asks a crowd
to decide which context should become visible \citep{wojcik2022birdwatch,
xcommunitynotes2026ranking}.

The key innovation of Community Notes is its bridging algorithm. A simple
majority vote would be easy to capture by one political side or one active
group. X's current 2026 official guide describes the main ranking layer as
matrix factorization over the sparse note-rater matrix. In simplified form, the
predicted rating is
\[
    \hat{r}_{un} = \mu + i_u + i_n + f_u \cdot f_n.
\]
This decomposes a rating into four parts:

\begin{itemize}
    \item $\mu$: the global average rating level;
    \item $i_u$: the rater intercept, capturing whether a rater is generally
    more generous or stricter when rating notes;
    \item $i_n$: the note intercept, used as the note's helpfulness score;
    \item $f_u \cdot f_n$: the match between the rater factor and the note
    factor, capturing whether a note appeals to raters with particular
    viewpoints.
\end{itemize}

The important point is that a note receives a high helpfulness score only when
its support cannot be explained merely by viewpoint compatibility. X therefore
presents the model as a way to identify notes with broad appeal across
viewpoints \citep{xcommunitynotes2026ranking}.

This matrix-factorization layer is not the whole current production system.
The official ranking guide also describes agreement safeguards, Net Helpful
minimums, checks for rater engagement and rater independence, topic and group
models, and a November 2025 Gaussian final-scoring pilot that approximates a
geometric mean of helpfulness rates across the rater-factor spectrum
\citep{xcommunitynotes2026ranking}. We therefore treat the formula as the core
ranking layer, not as the full Community Notes pipeline.

This design is powerful, but it also has a problem. After viewpoint agreement
is modeled, the system still estimates note helpfulness from sparse and uneven
ratings. Recent work argues that this makes the quality estimate sensitive to
rater quality, noisy raters, strategic raters, and hyperactive minorities
\citep{goyal2026qsmf,nudo2026hyperactive}. The official guide itself reflects
this concern: its rater-engagement safeguard recomputes a baseline intercept
after omitting the top 0.1\% of rating-volume contributors, in order to check
whether notes are helpful beyond the most active raters
\citep{xcommunitynotes2026ranking}. Participation in Community Notes is also
highly unequal. Razuvayevskaya et al. report that the top 10\% of contributors
produce 58\% of all notes, while most users remain much less active
\citep{razuvayevskaya2025timeliness}. This matters because a crowd system can
formally include many people while still depending heavily on a small, active
core.

The reason is structural. Community Notes faces sparse and uneven
participation: most raters see only a small part of the note universe, while a
small number of users rate very often. The current model answers this problem
as an estimation task. It adjusts for rater-level tendencies, estimates note
helpfulness, and uses a latent rater-note factor to capture whether a note fits
the raters who evaluate it. This is a reasonable way to prevent one-sided
agreement from being mistaken for broad quality. But it still keeps group
support inside a fitted model. It does not create direct quantities such as
``approval inside constituency A'' or ``approval inside constituency B.'' As a
result, cross-group acceptance is present only indirectly, and a small active
core can still shape which signals the model sees most often.

Our approach starts from a different premise about the disagreement itself.
Latent-quality estimation makes sense only if rater disagreement is noise around
a shared quality signal, something that washes out once it is partialled away.
On a platform that works as a de facto public sphere, that picture does not
hold. Raters do not fluctuate randomly around one true score; they cluster by
what they attend to: which claims are worth checking, what counts as credible
evidence, what context feels necessary. Those judgments cohere within a
group and pull apart across groups. The disagreement that results is neither
noise nor reducible to a single ideological axis. It is structured: persistent,
correlated across issues, and organized into behavioral blocs that we later
recover from real co-rating patterns.

Because the disagreement is structured, we do not try to decide which raters are
better, more sincere, or more reliable as individuals. We treat the recovered
constituencies as legitimate parts of the crowd and ask whether a note is
acceptable across them, which turns Community Notes from a prediction problem
into a crowd-agreement problem: the platform is making a collective decision, in
the name of ``the community,'' about whether a note should be attached to a
public post. The question this paper asks follows directly:

\begin{quote}
When ratings come from behaviorally distinct constituencies, which aggregation
rule should decide whether a Community Note becomes visible?
\end{quote}

We compare the platform's current decision with a rule that requires meaningful
support from each recovered constituency, treating disagreement as a feature of
the governance problem rather than as noise to be averaged away. This paper makes
three contributions:

\begin{itemize}[nosep,leftmargin=1.2em]
    \item We reframe the Community Notes display decision as an aggregation
    problem over behaviorally distinct constituencies, rather than a
    latent-quality estimate over a pooled crowd.
    \item We recover those constituencies endogenously from co-rating
    behavior, with no external party, country, or ideology labels, and give a
    transparent standing test for when a recovered cluster may count as a
    constituency.
    \item We operationalize cross-group legitimacy as a concrete
    non-compensatory rule (a geometric mean under minimum coverage) and show, on
    a 100{,}000-note slice, that it surfaces a large pool of otherwise-hidden
    public-interest notes, which a second-stage screen validates as sourced and
    factual.
\end{itemize}

The inspiration is political science and social choice rather than another
rater-quality model. Divided societies often use power sharing, double
majorities, parallel consent, vote pooling, or veto-like rules so that
collective decisions depend on support traveling across groups, not only on the
strongest or most active side
\citep{lijphart1977democracy,mcgarryoleary2009,
horowitz1991makingmoderation,reilly2002electoral}. Social-choice theory gives
a formal version of the same idea: legitimate collective decisions should
represent relevant parts of the public, not only the largest or most active
part \citep{peters2019proportionality,qiu2025representative}. We borrow this
logic in a limited way. Community Notes is not a state, and displaying a note
is not passing a law. Still, both settings face the same aggregation problem
under disagreement.

\begin{tcolorbox}[title=Philosophy, fonttitle=\bfseries,
    coltitle=black, colbacktitle=green!20, colback=green!6,
    colframe=green!45!black, boxrule=0.6pt, arc=2pt,
    left=6pt, right=6pt, top=5pt, bottom=5pt]
\textbf{Our philosophy is simple:} all recovered constituencies count. A note
should become visible when it can survive cross-group acceptance, not when one
active group dominates the observed rating record. Over time, such a rule may
also encourage contributors to write and rate notes that can travel across
camps, because visibility depends on cross-constituency support rather than
one-sided enthusiasm.
\end{tcolorbox}

\begin{tcolorbox}[title=Principles, fonttitle=\bfseries,
    coltitle=black, colbacktitle=green!20, colback=green!6,
    colframe=green!45!black, boxrule=0.6pt, arc=2pt,
    left=6pt, right=6pt, top=5pt, bottom=5pt]
\begin{description}[leftmargin=0pt,labelsep=0.6em,itemsep=0.25em]
    \item[\textbf{\emph{P1. Presence.}}] Every constituency must be present in the
    decision.
    \item[\textbf{\emph{P2. Non-compensation.}}] One constituency's rejection should
    not be fully compensated by another constituency's enthusiasm.
    \item[\textbf{\emph{P3. Symmetry.}}] Constituencies should enter the rule
    symmetrically, so that no group is privileged by default.
    \item[\textbf{\emph{P4. Behavioral recovery.}}] Constituencies should be
    recovered from observed rating behavior rather than assigned from outside
    labels such as party, country, or ideology.
\end{description}
\end{tcolorbox}

\medskip

We operationalize these rules as cross-constituency aggregation. Using the
active 200k Representative pipeline, we cluster active raters from co-rating
behavior, reassign them into two main behavioral constituencies, compute each
note's approval rate inside each constituency, and combine these rates with a
geometric mean under a minimum-coverage requirement. The geometric mean creates
a soft veto: if either constituency strongly rejects a note, the final score
falls, even if the other constituency is highly supportive. We then use this
score to identify hidden Community Notes that deserve a second look, and we
validate those candidates separately for factuality, sourcing, relevance, and
usefulness.

Concretely, let $a_{ic}$ be note $i$'s approval rate inside constituency $c$.
Writing $H_{ic}$ and $N_{ic}$ for the helpful and total ratings the note receives
from constituency $c$,
\[
a_{ic} =
\begin{cases}
H_{ic}/N_{ic}, & N_{ic} > 0,\\
0.5, & N_{ic} = 0,
\end{cases}
\]
where the $0.5$ fallback marks abstention rather than endorsement: a constituency
that has never rated the note casts no vote. The cross-constituency score is the
geometric mean
\[
    C_i = \sqrt{a_{i1}\,a_{i2}},
\]
and a note is selected when $C_i \geq 0.5$ and each constituency contributes at
least three raters; when Community Notes does not already display such a note, it
enters the rescue pool. Because $C_i$ falls toward zero as either approval rate
does, strong support from one constituency cannot rescue a note the other
rejects (P2), while the rule reads only within-constituency rates, so a larger or
busier bloc earns no extra weight (P3). The minimum-coverage requirement enforces
presence (P1), and the constituencies themselves are recovered from co-rating
behavior rather than assigned from outside labels (P4).

In a 100,000-note analytical slice, Community Notes displays 6,832 of 44,722
representative-picked notes, or 15\%. Cross-constituency aggregation identifies
a 13,655-note rescue pool, raising potential visibility to 20,487 notes, or
46\%, before validation. A second-stage validation screen classifies 3,896 of
these candidates as sourced, factual, and rescue-worthy. These results suggest
that Community Notes may have an aggregation problem as well as a rating
problem: useful public-interest notes can remain hidden when cross-group
support is modeled only indirectly.

\bibliographystyle{plainnat}
\bibliography{references}

\end{document}
